\documentclass{article}
\usepackage{amsmath,amssymb,amsthm}
\usepackage{algorithm}
\usepackage{algpseudocode}
\usepackage{hyperref}
\usepackage{enumitem}
\newtheorem{theorem}{Theorem}
\newtheorem{lemma}{Lemma}
\newtheorem{assumption}{Assumption}
\newtheorem{definition}{Definition}
\begin{document}

\section*{Cyclic Stochastic Gradient Descent (Cyclic SGD)}

\section*{Mathematical Formulation}
Let $f:\mathbb{R}^d\rightarrow\mathbb{R}$ be the (possibly non‑convex) objective and let 
$g_t$ be an unbiased stochastic gradient computed on a random mini‑batch at iteration $t$:
\[
\mathbb{E}[g_t\mid w_t]=\nabla f(w_t),\qquad 
\mathbb{E}\big[\|g_t-\nabla f(w_t)\|^2\mid w_t\big]\le\sigma^2 .
\]

A \emph{triangular cyclic learning‑rate schedule} with period $2M$ is defined by
\[
\eta_t \;:=\; \eta_{\min}+(\eta_{\max}-\eta_{\min})\;
\Bigl(1-\Bigl|\frac{(t-1)\bmod (2M)}{M}-1\Bigr|\Bigr),\qquad t=1,2,\dots
\tag{1}
\]
Thus $\eta_t$ linearly increases from $\eta_{\min}$ to $\eta_{\max}$ in the first $M$ steps,
then linearly decreases back to $\eta_{\min}$ in the next $M$ steps, and the pattern repeats.

The update rule is the standard SGD step with the time‑varying step size:
\[
w_{t+1}=w_t-\eta_t\,g_t,\qquad t=0,1,2,\dots
\tag{2}
\]

Hyper‑parameters:
\begin{itemize}
    \item $\eta_{\min}>0$, $\eta_{\max}> \eta_{\min}$ – minimal and maximal step sizes,
    \item $M\in\mathbb{N}$ – half‑period of the cycle,
    \item $T$ – total number of iterations (multiple of $2M$ for simplicity).
\end{itemize}

\section*{Pseudocode}
\begin{algorithm}[H]
\caption{Cyclic SGD}
\begin{algorithmic}[1]
\Require Initial point $w_0\in\mathbb{R}^d$, 
        $\eta_{\min},\eta_{\max}>0$, period $M$, total steps $T$.
\For{$t=0$ \textbf{to} $T-1$}
    \State Compute stochastic gradient $g_t$ on a random mini‑batch.
    \State Compute cyclic step size
    \[
    \eta_t \gets \eta_{\min}+(\eta_{\max}-\eta_{\min})
    \Bigl(1-\Bigl|\frac{t\bmod (2M)}{M}-1\Bigr|\Bigr)
    \tag{3}
    \]
    \State Update $w_{t+1}\gets w_t-\eta_t g_t$.
\EndFor
\State \textbf{return} $w_T$ (or optionally the average $\bar w_T=\frac1T\sum_{t=0}^{T-1}w_t$).
\end{algorithmic}
\end{algorithm}

\section*{Dry Run Example}
Consider the one‑dimensional quadratic $f(w)=(w-3)^2$.
\[
\nabla f(w)=2(w-3).
\]
Choose $w_0=0$, $\eta_{\min}=0.05$, $\eta_{\max}=0.1$, $M=2$ (period $4$).  
We perform three iterations.

\begin{center}
\begin{tabular}{c|c|c|c|c}
$t$ & $\eta_t$ (from (3)) & $g_t=\nabla f(w_t)$ & $w_{t+1}=w_t-\eta_t g_t$ & $f(w_{t+1})$ \\ \hline
0 & $0.05$ & $2(0-3)=-6$ & $0-0.05(-6)=0.30$ & $(0.30-3)^2=7.29$ \\
1 & $0.075$ & $2(0.30-3)=-5.40$ & $0.30-0.075(-5.40)=0.705$ & $(0.705-3)^2=5.27$ \\
2 & $0.10$ & $2(0.705-3)=-4.59$ & $0.705-0.10(-4.59)=1.164$ & $(1.164-3)^2=3.37$
\end{tabular}
\end{center}
The objective value decreases at each step, illustrating the effect of the cyclic step size.

\newpage

\section*{Assumptions}
\begin{assumption}[Smoothness]\label{ass:smooth}
$f$ is $L$‑smooth, i.e.,
\[
f(y)\le f(x)+\langle\nabla f(x),y-x\rangle+\frac{L}{2}\|y-x\|^2,
\qquad\forall x,y\in\mathbb{R}^d .
\tag{4}
\]
\end{assumption}

\begin{assumption}[Unbiased Gradient and Bounded Variance]\label{ass:sgd}
For every $t$,
\[
\mathbb{E}[g_t\mid w_t]=\nabla f(w_t),\qquad
\mathbb{E}\big[\|g_t-\nabla f(w_t)\|^2\mid w_t\big]\le\sigma^2 .
\tag{5}
\]
\end{assumption}

\begin{assumption}[Step‑Size Bounds]\label{ass:steps}
The cyclic schedule satisfies
\[
0<\eta_{\min}\le\eta_t\le\eta_{\max},\qquad\forall t,
\quad\text{and}\quad \eta_{\max}\le\frac{1}{L}.
\tag{6}
\]
\end{assumption}

\section*{Main Convergence Theorem}
\begin{theorem}[Convergence of Cyclic SGD]\label{thm:main}
Under Assumptions~\ref{ass:smooth}–\ref{ass:steps} and for any integer $T\ge1$,
\[
\frac{1}{T}\sum_{t=0}^{T-1}\mathbb{E}\big[\|\nabla f(w_t)\|^2\big]
\;\le\;
\frac{2\big(f(w_0)-f^\star\big)}{\eta_{\min}T}
\;+\;
L\,\eta_{\max}\,\sigma^2 .
\tag{7}
\]
Consequently, if $T\to\infty$, the average squared gradient norm converges to
$O\!\big(\frac{1}{T}\big)+O\!\big(\eta_{\max}\big)$. Choosing $\eta_{\max}=O(1/\sqrt{T})$ yields the classic non‑convex rate
$\mathbb{E}\big[\|\nabla f(\bar w_T)\|^2\big]=O(1/\sqrt{T})$.
\end{theorem}

\section*{Theoretical Properties}
\begin{enumerate}[label=\arabic*.]
    \item \textbf{Stability.} Because $\eta_{\max}\le 1/L$, the descent lemma (Lemma~\ref{lem:descent}) guarantees that each iteration does not increase the expected function value by more than a term proportional to the variance.
    \item \textbf{Hyper‑parameter Sensitivity.} The bound (7) depends only on the extreme values $\eta_{\min}$ and $\eta_{\max}$, not on the exact shape of the cycle. Hence the method is robust to the choice of $M$ as long as the period is not too short to violate (6).
    \item \textbf{Comparison to Baselines.} Standard SGD with a constant step size $\eta$ satisfies the same bound with $\eta_{\min}=\eta_{\max}=\eta$. Cyclic SGD can exploit a larger $\eta_{\max}$ (for faster escape from saddles) while keeping $\eta_{\min}$ small enough to guarantee stability, often leading to better empirical performance.
\end{enumerate}

\section*{Discussion}
The theorem shows that the cyclic schedule does not deteriorate the worst‑case convergence guarantee of SGD. Moreover, the presence of a larger $\eta_{\max}$ can improve practical speed of descent, especially in regions where the gradient norm is small. The bound is tight up to constants: if $\sigma=0$ (full‑batch case) the term $L\eta_{\max}\sigma^2$ vanishes and (7) reduces to the deterministic gradient descent rate.

\newpage

\section*{Preliminaries (Definitions and Lemmas)}
\begin{definition}[Filtration]
Let $\mathcal{F}_t:=\sigma\{w_0,g_0,\dots,g_{t-1}\}$ be the $\sigma$‑algebra generated by all randomness up to iteration $t$.
\end{definition}

\begin{lemma}[Descent Lemma for $L$‑smooth functions]\label{lem:descent}
For any $w\in\mathbb{R}^d$ and any vector $d$,
\[
f(w+d)\le f(w)+\langle\nabla f(w),d\rangle+\frac{L}{2}\|d\|^2 .
\tag{8}
\]
\end{lemma}
\begin{proof}
Directly from Assumption~\ref{ass:smooth} with $x=w$, $y=w+d$.
\end{proof}

\section*{Main Proof (Step‑by‑Step Derivation)}
We prove Theorem~\ref{thm:main}. Throughout we condition on $\mathcal{F}_t$ and use the notation
$\mathbb{E}_t[\cdot]=\mathbb{E}[\cdot\mid\mathcal{F}_t]$.

\paragraph{Step 1. Apply the descent lemma to the SGD update.}
From (8) with $d=-\eta_t g_t$,
\[
f(w_{t+1})\le f(w_t)-\eta_t\langle\nabla f(w_t),g_t\rangle
+\frac{L}{2}\eta_t^{2}\|g_t\|^{2}.
\tag{9}
\]
[Step 1]

\paragraph{Step 2. Take conditional expectation.}
Using unbiasedness (5) and the fact that $\eta_t$ is $\mathcal{F}_t$‑measurable,
\[
\mathbb{E}_t\!\big[f(w_{t+1})\big]
\le f(w_t)-\eta_t\|\nabla f(w_t)\|^{2}
+\frac{L}{2}\eta_t^{2}\,\mathbb{E}_t\!\big[\|g_t\|^{2}\big].
\tag{10}
\]
[Step 2]

\paragraph{Step 3. Expand $\mathbb{E}_t[\|g_t\|^{2}]$.}
\[
\mathbb{E}_t\!\big[\|g_t\|^{2}\big]
= \|\nabla f(w_t)\|^{2}
   +\mathbb{E}_t\!\big[\|g_t-\nabla f(w_t)\|^{2}\big]
\le \|\nabla f(w_t)\|^{2}+\sigma^{2},
\tag{11}
\]
where the inequality follows from variance bound (5).  
[Step 3]

\paragraph{Step 4. Substitute (11) into (10).}
\[
\mathbb{E}_t\!\big[f(w_{t+1})\big]
\le f(w_t)-\eta_t\|\nabla f(w_t)\|^{2}
+\frac{L}{2}\eta_t^{2}\Big(\|\nabla f(w_t)\|^{2}+\sigma^{2}\Big).
\tag{12}
\]
[Step 4]

\paragraph{Step 5. Rearrange terms to isolate $\|\nabla f(w_t)\|^{2}$.}
\[
\mathbb{E}_t\!\big[f(w_{t+1})\big]
\le f(w_t)-\Big(\eta_t-\frac{L}{2}\eta_t^{2}\Big)\|\nabla f(w_t)\|^{2}
+\frac{L}{2}\eta_t^{2}\sigma^{2}.
\tag{13}
\]
[Step 5]

\paragraph{Step 6. Use the step‑size bound $\eta_{\max}\le 1/L$.}
Since $0<\eta_t\le\eta_{\max}\le 1/L$, we have
\[
\eta_t-\frac{L}{2}\eta_t^{2}
\;\ge\;
\frac{\eta_t}{2}
\;\ge\;
\frac{\eta_{\min}}{2}.
\tag{14}
\]
[Step 6]

\paragraph{Step 7. Apply (14) to (13).}
\[
\mathbb{E}_t\!\big[f(w_{t+1})\big]
\le f(w_t)-\frac{\eta_{\min}}{2}\|\nabla f(w_t)\|^{2}
+\frac{L}{2}\eta_{\max}^{2}\sigma^{2}.
\tag{15}
\]
[Step 7]

\paragraph{Step 8. Take full expectation and sum over $t=0,\dots,T-1$.}
Define $f^\star:=\inf_{w}f(w)$. Summing (15) and telescoping yields
\[
\frac{\eta_{\min}}{2}\sum_{t=0}^{T-1}
\mathbb{E}\!\big[\|\nabla f(w_t)\|^{2}\big]
\le f(w_0)-\mathbb{E}\!\big[f(w_T)\big]
+ \frac{L}{2}\eta_{\max}^{2}\sigma^{2}T.
\tag{16}
\]
Since $f(w_T)\ge f^\star$, we obtain
\[
\frac{1}{T}\sum_{t=0}^{T-1}
\mathbb{E}\!\big[\|\nabla f(w_t)\|^{2}\big]
\le
\frac{2\big(f(w_0)-f^\star\big)}{\eta_{\min}T}
+L\,\eta_{\max}\,\sigma^{2}.
\tag{17}
\]
[Step 8]

Equation (17) is exactly the bound (7) claimed in Theorem~\ref{thm:main}. ∎

\section*{Rate Derivation (With Constants)}
From (17) we read off the two contributions:
\begin{itemize}
    \item \emph{Optimization term} $ \displaystyle \frac{2\big(f(w_0)-f^\star\big)}{\eta_{\min}T}$ decays as $1/T$.
    \item \emph{Stochastic term} $L\eta_{\max}\sigma^{2}$ is constant for a fixed $\eta_{\max}$.
\end{itemize}
If we let the maximal step size decay as $\eta_{\max}=c/\sqrt{T}$ (with $c\le 1/\sqrt{L}$ to preserve (6)), then
\[
L\eta_{\max}\sigma^{2}= \frac{Lc\sigma^{2}}{\sqrt{T}} = O\!\big(\tfrac{1}{\sqrt{T}}\big).
\]
Thus the overall rate becomes
\[
\frac{1}{T}\sum_{t=0}^{T-1}\mathbb{E}\!\big[\|\nabla f(w_t)\|^{2}\big]
= O\!\Bigl(\frac{1}{T}\Bigr)+O\!\Bigl(\frac{1}{\sqrt{T}}\Bigr)
= O\!\Bigl(\frac{1}{\sqrt{T}}\Bigr).
\]
This matches the optimal non‑convex SGD rate.

\section*{Special Cases}
\begin{enumerate}[label=\alph*.]
    \item \textbf{Deterministic case ($\sigma=0$).} The bound reduces to
    \[
    \frac{1}{T}\sum_{t=0}^{T-1}\|\nabla f(w_t)\|^{2}
    \le \frac{2\big(f(w_0)-f^\star\big)}{\eta_{\min}T},
    \]
    i.e., a pure $O(1/T)$ decay.
    \item \textbf{Constant step size.} Setting $\eta_{\min}=\eta_{\max}=\eta$ recovers the classic SGD bound
    \[
    \frac{1}{T}\sum_{t=0}^{T-1}\mathbb{E}\!\big[\|\nabla f(w_t)\|^{2}\big]
    \le \frac{2\big(f(w_0)-f^\star\big)}{\eta T}+L\eta\sigma^{2}.
    \]
    \item \textbf{Very short cycles ($M=1$).} The schedule still satisfies (6) and the same proof applies; the bound is independent of $M$, confirming that the cycle length does not affect the worst‑case guarantee.
\end{enumerate}

\end{document}